\section{Architecture and Primitives}
\label{sec:receipts-architecture}

\subsection{Core Objects}
\label{sec:receipts-architecture-objects}

The receipts corpus is built on six core cryptographic and type-theoretic objects. The
first is the \textbf{Ed25519 validator signature}, carried in the \texttt{validator\_signature}
field of every M\&A claim receipt. This signature is computed over the JCS canonical form
(RFC 8785) of the JSON receipt, ensuring that the signature covers a deterministic byte
sequence regardless of key ordering in the source document. Any modification to any field
in the JSON receipt after signing invalidates the signature, making post-hoc tampering
detectable without access to the private key.

The second core object is the \textbf{BLAKE3 hash}, appearing as both \texttt{target\_log\_hash}
and \texttt{process\_model\_hash} in each M\&A receipt. BLAKE3 is a cryptographic hash
function that provides content addressing: two event logs with the same hash are guaranteed
to have the same content. This binding ensures that the conformance score reported in a receipt
was computed over the specific event log identified by the hash, not a different or modified log.

The third core object is the \textbf{ConformanceVerdicts} payload type, a Rust struct carrying
four quality metrics: fitness, precision, generalization, and simplicity. These four metrics
correspond to the van der Aalst quality framework for process-mining results. The
\texttt{ConformanceAdmissionGate} (Blue River Dam Gate 3) evaluates the fitness and precision
fields of this struct against thresholds $\geq 0.95$ and $\geq 0.90$ respectively, and refuses
to emit a board-admissible receipt if either threshold is not met.

The fourth core object is the \textbf{Evidence\textless{}T, State, Witness\textgreater{}} Rust
generic type, which enforces lifecycle state at the type-law boundary. This type ensures that
evidence objects can only transition through lawful states---from raw to validated to admitted---
and that the witness type parameter records which entity authorized each transition. The
fifth core object is the \textbf{RECEIPT\_REGISTRY}, a governance ledger enumerating seven
canonical receipt types. The sixth is the \textbf{JCS canonical serialization} layer (RFC 8785)
that makes JSON hashing deterministic.

\subsection{Admissible Transitions}
\label{sec:receipts-architecture-transitions}

The receipt lifecycle enforces three admissible transitions. The first is \emph{manufacture},
in which a wasm4pm query is executed over an event log and the resulting
\texttt{ConformanceVerdicts} payload is assembled into a JSON receipt document. At this stage
the receipt exists but carries no signature. The second transition is \emph{gate evaluation},
in which the \texttt{ConformanceAdmissionGate} evaluates the fitness and precision fields against
the board admissibility thresholds. If the gate passes, the receipt is forwarded for signing.
If the gate fails (fitness below $0.95$), the receipt is rejected and no signature is emitted.

The third transition is \emph{sealing}, in which the Ed25519 signing authority computes the
JCS canonical form of the JSON document and signs it. The signed receipt is then eligible for
inclusion in the M\&A board record and for citation in doctoral research claims. The
\texttt{RECEIPT\_REGISTRY.md} governs which receipt types may enter the board record and
which witnesses must be named for each type.

The wasm4pm module generation receipts follow a parallel but distinct transition sequence:
\emph{compile}, \emph{test}, \emph{hash}, \emph{seal}. The Blake3 hash of the compiled
module binds the generation receipt to a specific binary artifact. The test pass count
(e.g., $8/8$ for conformance authority) is recorded in the receipt body as the test witness.

\subsection{Boundaries}
\label{sec:receipts-architecture-boundaries}

The receipts corpus enforces four architectural boundaries. The first is the
\textbf{manufacture-to-claim boundary}: no financial claim may enter the M\&A record without
a receipted conformance measurement. The \texttt{board\_admissibility} field in each JSON
receipt is the machine-readable assertion that this boundary has been crossed lawfully.

The second is the \textbf{module-to-output boundary}: no wasm4pm module output may be used
in a claim unless the module's generation receipt exists and documents a passing compilation
and test suite. The conformance authority generation receipt for version \texttt{v30.1.2}
documents this boundary with Blake3 hash
\texttt{75508e6e9e7acccb679e5c8be35ddd2adc9c4732f1fc331047d755703c9107c3}.

The third is the \textbf{research-to-execution boundary}: no downstream wasm4pm refactor
may proceed until the research program's receipt chain has spoken. This boundary is
enforced by doctrine, not by code, and is documented in the process-intelligence
\texttt{CLAUDE.md} configuration.

The fourth is the \textbf{assertion-to-evidence boundary}: no receipt type may carry a
claim derived from inference alone. Claims derived from assumption or inference are
\texttt{PARTIAL} findings, documented in \texttt{gaps/} rather than the signed receipt chain.

\subsection{Failure Modes Refused}
\label{sec:receipts-architecture-failures}

The receipts architecture refuses four specific failure modes by design. The first is
\emph{fitness inflation}: the \texttt{ConformanceAdmissionGate} hard-rejects any receipt
with fitness below $0.95$, preventing a low-conformance result from being represented
as board-admissible. The unit test \texttt{test\_admission\_gate\_hard\_reject} verifies
this behavior programmatically.

The second refused failure mode is \emph{log substitution}: because each receipt carries
a BLAKE3 \texttt{target\_log\_hash}, substituting a different event log after signing
invalidates the hash binding and can be detected by recomputing the hash of the claimed
log.

The third refused failure mode is \emph{signature forgery}: the Ed25519 scheme is
existentially unforgeable under chosen-message attack, meaning no party without the
private key can manufacture a valid signature for a modified receipt.

The fourth refused failure mode is \emph{premature board admission}: the
\texttt{board\_admissibility} assessment object is only populated after the
\texttt{ConformanceAdmissionGate} passes, ensuring that a receipt in an intermediate
manufacturing state cannot be mistaken for a finalized, board-admissible claim.
